\documentclass[11pt,a4paper]{article}

\usepackage[utf8]{inputenc}
\usepackage[T1]{fontenc}
\usepackage{lmodern}
\usepackage[english]{babel}
\usepackage{microtype}
\usepackage[margin=1in]{geometry}
\usepackage{amsmath,amssymb}
\usepackage{booktabs}
\usepackage{array}
\usepackage{graphicx}
\usepackage{xcolor}
\usepackage{url}
\usepackage{hyperref}
\usepackage{cleveref}
\usepackage[style=numeric,backend=biber,maxbibnames=99,maxcitenames=2,
            url=false,doi=true]{biblatex}
\addbibresource{references.bib}

\hypersetup{
  colorlinks=true,
  citecolor=blue!60!black,
  linkcolor=blue!60!black,
  urlcolor=blue!60!black,
  pdftitle={KVScope V2: Cross-Architecture KV-Cache Profiling on NVIDIA H200},
  pdfauthor={Rahul Surya}
}

\graphicspath{{figures/}}

\newcommand{\code}[1]{\texttt{#1}}
\newcommand{\NA}{--}

\title{\textbf{KVScope V2: Cross-Architecture KV-Cache Profiling\\
       on NVIDIA H200}}
\author{Rahul Surya\\[4pt]
  \small MSc HPC with Data Science, EPCC, University of Edinburgh\\
  \small ORCID: 0009-0009-8629-7776}
\date{May 2026}

\begin{document}
\maketitle

% =============================================================
\begin{abstract}
Key--value (KV) caches account for the dominant growth term in
autoregressive inference memory, yet per-architecture footprint
fingerprints remain absent from most model cards.
We apply KVScope~V2 to six model families spanning four
attention architectures on an NVIDIA~H200
(143\,771\,MiB VRAM, Hopper architecture).
All models are profiled under a 2048-token generation budget
(1024 for the smaller Pythia-1.4B baseline), substantially
reducing the token-budget confound of prior work.
Three concrete results emerge:
(1)~KV-cache density is a stable, architecture-determined
fingerprint---dense MHA measures 8.00\,KB/token/layer,
dense GQA measures 4.00\,KB/token/layer, and hybrid
SSM+attention models range from 1.12 to 2.00\,KB/token/layer;
(2)~Gemma~4's local/global hybrid produces the highest density
at 14.67\,KB/token/layer with a coefficient of variation of
0.203, confirming bimodal per-layer footprints;
(3)~gpt-oss exhibits extreme layer heterogeneity
(CV\,=\,0.902) from alternating sliding-window and
full-attention layers, the highest CV in the KVScope dataset.
Nemotron-H is retained as a state-space-model control with
\code{use\_cache=False}, confirming O(1) memory trajectory.
Speculative decoding with a 70B verifier and 1B draft model
shows 1.8\,GB dual-model NVML overhead at 2048 tokens.
\end{abstract}

\noindent\textbf{Keywords:} KV cache; transformer inference;
NVIDIA H200; grouped-query attention; sliding-window attention;
state-space models; speculative decoding; memory profiling.

% =============================================================
\section{Introduction}
\label{sec:intro}

Autoregressive decoding reuses previously computed keys and values
to avoid recomputing attention over the full prefix at each step.
This optimization converts compute into persistent memory pressure:
each active attention layer retains a KV tensor whose size scales
with sequence length, number of KV heads, head dimension, and
element width.
Despite this predictable growth, the operational memory footprint
of individual models is rarely disclosed: model cards report
parameter counts and quantisation precision, but not the KB of KV
cache consumed per generated token per layer.
The observability gap matters because serving systems allocate
GPU memory in advance; miscalibrated estimates lead either to
out-of-memory crashes or to wasteful over-provisioning
\cite{kwon2023pagedattention}.

The original KVScope study addressed this gap for six model
families on an NVIDIA~H100 and published a reproducible profiling
framework~\cite{kvscope2026h100}.
A follow-up on the H200~(V1) extended coverage but suffered from
inconsistent token budgets (300--4096 tokens across artifacts),
making cross-model comparisons unreliable.
This paper reports KVScope~V2, which re-profiles all models under
a 2048-token generation budget (1024 for Pythia-1.4B) on the same
H200 hardware.
The goals are:
(a)~to establish architecture-level KV-cache density as a stable,
GPU-independent fingerprint;
(b)~to quantify layer heterogeneity in hybrid architectures;
(c)~to measure the KV overhead of speculative decoding; and
(d)~to document practical profiling failures that matter for
reproducible systems research.
We do \emph{not} claim improved throughput over the H100 study
because protocols and run quality differ across artifacts.

% =============================================================
\section{Background}
\label{sec:background}

\paragraph{Transformer KV-cache scaling.}
For a dense transformer with batch size~$B$, sequence
length~$S$, number of KV heads~$H_\mathrm{kv}$, head
dimension~$D$, and element size~$E$\,bytes, the stored KV cache
per layer is
\begin{equation}
  M_\mathrm{KV,layer} = 2\,B\,S\,H_\mathrm{kv}\,D\,E,
  \label{eq:kvcache}
\end{equation}
where the factor of two accounts for keys and values separately.
Multi-query attention sets $H_\mathrm{kv}=1$; grouped-query
attention (GQA) uses $H_\mathrm{kv}<H_\mathrm{q}$~\cite{shazeer2019mqa,ainslie2023gqa}.
Sliding-window layers cap $S$ at a fixed window, producing a
bimodal per-layer footprint in hybrid architectures
\cite{beltagy2020longformer}.
Paged serving further complicates accounting because
\texttt{torch.cuda.memory\_reserved} can diverge from live tensor
bytes~\cite{kwon2023pagedattention}.

\paragraph{Hybrid SSM/attention architectures.}
Modern large models increasingly replace some attention layers with
state-space model (SSM) blocks~\cite{gu2023mamba} or linear
attention variants.
These blocks maintain a fixed-size recurrent state rather than a
growing KV tensor; memory is thus O(1) with respect to sequence
length rather than O($S$).
Qwen~3.6 uses DeltaNet layers that bypass KV storage entirely,
leaving only 10 of 40 layers with active KV tensors.
gpt-oss alternates sliding-window and full-attention layers,
producing extreme per-layer heterogeneity.
Gemma~4 uses a local/global hybrid where 50 sliding-attention
layers use sliding windows and 10 full-attention layers attend to
the full sequence.
KVScope registers hooks only on standard attention modules and
therefore captures no KV snapshots for SSM or DeltaNet layers,
making such models useful as zero-KV-cache controls or
partial-coverage profiling targets.

\paragraph{Speculative decoding.}
Speculative decoding uses a small draft model to propose $k$ tokens
which a larger verifier model validates in parallel~\cite{leviathan2023fast}.
The KV overhead is amplified because both models maintain
independent KV caches during the draft--verify cycle.
At sequence length $S$, the combined KV footprint is approximately
\begin{equation}
  M_\mathrm{spec} \approx M_\mathrm{KV}^{(\mathrm{verifier})}(S)
    + k \cdot M_\mathrm{KV}^{(\mathrm{draft})}(S),
  \label{eq:spec}
\end{equation}
where $k$ is the number of draft tokens per step.

% =============================================================
\section{Methods}
\label{sec:methods}

\subsection{Hook-based KV extraction}
\label{sec:methods:hooks}

The profiler registers a \code{KVSnapshot} forward hook on every
attention module identified by the model adapter.
At each step the hook records: layer index, layer type
(\emph{local}/\emph{global}/\emph{dense}), K and V
tensor shapes and byte counts, dtype, a timestamp, and PyTorch
allocator counters (\code{torch.cuda.memory\_reserved},
\code{torch.cuda.memory\_allocated}).
A separate NVML poller samples GPU memory every few steps via
\texttt{pynvml}.

The final per-layer byte count for each prompt is the last hook
snapshot before the EOS token.
Total KV footprint is the sum over all profiled layers.
The cache-density descriptor divides each layer's final byte count
by its observed sequence length:
\begin{equation}
  \rho_l = \frac{M_\mathrm{KV,layer}^{(l)}}{S_l}
           \quad\text{(bytes per token for layer }l\text{)},
\end{equation}
and the per-model density is the median $\rho_l$ across all active
layers and prompts.

\subsection{Detector pipeline}
\label{sec:methods:detectors}

Five detectors run per prompt and their scores are aggregated by a
fixed weighted sum:
\begin{equation}
  s_\mathrm{total}
  = 0.35\,s_\mathrm{postEOS}
  + 0.30\,s_\mathrm{growth}
  + 0.20\,s_\mathrm{frag}
  + 0.10\,s_\mathrm{layer}
  + 0.05\,s_\mathrm{mla}.
  \label{eq:score}
\end{equation}

\paragraph{Post-EOS detector.}
Let $M_\mathrm{base}$, $M_\mathrm{peak}$, and $M_\mathrm{post}$
denote NVML samples at generation start, maximum, and after EOS.
\begin{align}
  \Delta_\mathrm{freed}
    &= M_\mathrm{peak} - M_\mathrm{base}, \label{eq:freed}\\
  \mathrm{retained}_{MB}
    &= M_\mathrm{post} - M_\mathrm{base},  \label{eq:retained}\\
  s_\mathrm{postEOS}
    &= \mathrm{clip}\!\left(
         \frac{\mathrm{retained}_{MB}}{\Delta_\mathrm{freed}},
         0, 1
       \right). \label{eq:posteos}
\end{align}

\paragraph{Fragmentation detector.}
\begin{equation}
  \mathrm{frag}_{MB}
  = \texttt{cuda.memory\_reserved} - \texttt{cuda.memory\_allocated}.
  \label{eq:frag}
\end{equation}

\paragraph{Growth detector.}
Fits linear and quadratic curves to the total KV footprint over
decode steps.
A high score requires both a substantial fit improvement for the
quadratic term and an upward residual trend from the linear fit.
A plateau guard suppresses false positives on sliding-window
architectures.

\paragraph{Layer-anomaly detector.}
Identifies per-layer cache outliers beyond three standard
deviations of the per-layer final distribution.

% =============================================================
\section{Experimental Setup}
\label{sec:setup}

\subsection{Hardware and software}
\label{sec:setup:hw}

All runs targeted a single NVIDIA~H200 SXM5 GPU
(143\,771\,MiB VRAM, Hopper architecture, driver~590.48.01,
PCIe~Gen5$\times$16).
The host system runs Ubuntu~22.04.5 LTS on an Intel Xeon
Platinum~8592+ (24~cores at 1.9\,GHz).
Software versions: Python~3.10.12, PyTorch~2.12.0+cu130,
CUDA~13.0, Triton~3.7.0, Transformers~5.9.0,
bitsandbytes~0.49.2, accelerate~1.13.0.
The full \code{pip freeze} and \code{nvidia-smi -q} output are
included in the artifact package.

\subsection{Models}
\label{sec:setup:models}

\begin{table*}[t]
  \centering
  \small
  \caption{Verified H200 V2 profiling metrics.
    Density and CV are computed from per-layer KV tensor byte
    counts captured by forward hooks.
    Nemotron-H has no KV snapshots (SSM architecture,
    \code{use\_cache=False}).
    Spec decode uses 4-bit NF4 quantisation for the 70B verifier.}
  \label{tab:main-metrics}
  \resizebox{\textwidth}{!}{%
  \begin{tabular}{@{}llrrrrrrrc@{}}
    \toprule
    \textbf{Model} & \textbf{Architecture} & \textbf{Tokens} & \textbf{Valid} & \textbf{Layers} & \textbf{KB/tok/lay} & \textbf{CV} & \textbf{KV MB} & \textbf{Leak max} & \textbf{Grade} \\
    \midrule
    Pythia-1.4B       & Dense MHA              & 1024 & 20/20 & 24 & 8.00  & 0.000 & 210.2  & 0.888 & A \\
    Gemma~4 (27B)     & Hybrid Local/Global    & 2048 & 20/20 & 60 & 14.67 & 0.203 & 1892.3 & 0.330 & A \\
    Qwen~3.6 (35B)    & Hybrid DeltaNet+GQA     & 2048 & 20/20 & 10 & 2.00  & 0.000 & 42.4   & 0.305 & A \\
    gpt-oss (Hybrid)  & Hybrid SSM+Attention   & 2048 & 20/20 & 36 & 1.12  & 0.902 & 79.9   & 0.222 & A \\
    Nemotron-H (SSM)  & Pure SSM (no KV)       & 2048 & 20/20 & \NA & \NA & \NA & \NA & 0.000 & D \\
    Spec Decode (70B+1B) & Speculative Decoding & 2048 &  8/20 & \NA & \NA & \NA & 1804.6\rlap{$^*$} & 0.000 & B \\
    \bottomrule
  \end{tabular}}
  \vspace{2pt}
  {\footnotesize $^*$NVML overhead (includes model weights for both models).}
\end{table*}

\Cref{tab:main-metrics} summarises the six model families.
Pythia-1.4B~\cite{biderman2023pythia} serves as the pure-MHA
baseline with 24~layers and 16~attention heads.
Gemma~4~\cite{gemma4_2026} uses a local/global hybrid with 50
sliding-attention and 10~full-attention layers across 60~total
attention layers; sliding layers use 16~KV heads with head
dimension~256, while full-attention layers use 4~KV heads with
head dimension~512.
Qwen~3.6~\cite{qwen2025qwen3} is a DeltaNet+GQA hybrid where
only 10 of 40~layers produce KV tensors; the remaining 30~DeltaNet
layers maintain fixed recurrent states.
gpt-oss~\cite{openai2025gptoss} alternates sliding-window and
full-attention layers across 36~profiled layers, with 8~KV heads
and head dimension~64.
Nemotron-H~\cite{nvidia2025nemotron} is a Mamba-SSM hybrid run
with \code{use\_cache=False} as an architectural control.
The speculative decoding configuration pairs Llama-3.1-70B
(80~layers, 8~KV heads, head dimension~128) as verifier with
Llama-3.2-1B (16~layers) as draft, using 4-bit NF4 quantisation
for the verifier to fit within H200 VRAM.

\subsection{Protocol}
\label{sec:setup:protocol}

All models are profiled with \code{MAX\_TOKENS=2048} and
\code{CAPTURE\_EVERY=5}, using 20~prompts drawn from
\texttt{experiments/configs/prompts\_advanced.json}.
The prompts include long-form generation tasks (code, essays,
stories), short-answer questions, and adversarial repeats to
exercise both full-length and early-termination paths.
This 2048-token budget (1024 for Pythia-1.4B, whose shorter
context window prevents longer generation) eliminates the
generation-length confound that plagued V1, where token budgets
ranged from 300 to 4096 across artifacts.

Compared to the H100 study~(V1) and the H200~V1 follow-up:
\begin{itemize}\setlength\itemsep{1pt}
  \item Token budget is 2048 for all models except Pythia
    (1024); V1 used 300--4096 tokens.
  \item Qwen~3.6 replaces the failed LFM2.5-350M profile
    (19/20 prompts terminated at 0--9 steps in V1 due to
    missing chat template).
  \item Speculative decoding is added as a new experiment.
  \item 4-bit NF4 quantisation is used for the 70B verifier
    in speculative decoding to fit within H200~VRAM.
  \item Environment metadata is fully captured: \code{pip freeze},
    \code{nvidia-smi -q}, and CPU/OS details are included in
    the artifact package.
\end{itemize}

% =============================================================
\section{Results}
\label{sec:results}

The analysis focuses on two primary metrics: cache density
(KB/token/layer) and per-layer coefficient of variation (CV).
The detector pipeline scores from \cref{eq:score} are recorded
in \cref{tab:main-metrics} for completeness; detailed detector
analysis is deferred to future work.

\subsection{Architecture-level cache density}
\label{sec:results:density}

The principal finding is that cache density follows architecture,
not GPU generation or model scale.
\Cref{tab:density-hierarchy} presents the density hierarchy.

\begin{table}[t]
  \centering
  \small
  \caption{KV-cache density hierarchy across architectures.
    All values in KB/token/layer, computed from per-layer
    KV tensor byte counts.}
  \label{tab:density-hierarchy}
  \begin{tabular}{@{}lrc@{}}
    \toprule
    \textbf{Architecture} & \textbf{KB/tok/layer} & \textbf{Model} \\
    \midrule
    Hybrid SSM+Attn (alternating) & 1.12 & gpt-oss \\
    Hybrid DeltaNet+GQA           & 2.00 & Qwen~3.6 \\
    Dense GQA                     & 4.00 & Llama-70B\rlap{$^\dagger$} \\
    Dense MHA                     & 8.00 & Pythia-1.4B \\
    Hybrid Local/Global           & 14.67 & Gemma~4 \\
    \bottomrule
  \end{tabular}
  \vspace{2pt}
  {\footnotesize $^\dagger$From speculative decoding verifier
  (8/20 valid prompts).}
\end{table}

The density hierarchy is consistent with theoretical predictions
from \cref{eq:kvcache}:
dense MHA with 16~KV heads and head dimension~128 produces
$2 \times 16 \times 128 \times 2 = 8192$~bytes/token/layer
(8.00\,KB);
GQA with 8~KV heads and head dimension~128 produces
$2 \times 8 \times 128 \times 2 = 4096$~bytes/token/layer
(4.00\,KB);
gpt-oss with 8~KV heads and head dimension~64 produces
$2 \times 8 \times 64 \times 2 = 2048$~bytes/token/layer
for full-attention layers, but the alternating sliding-window
layers reduce the effective density to 1.12\,KB/token/layer.

Gemma~4 exceeds the pure-MHA baseline because its 50~sliding
layers use head dimension~256 with 16~KV heads, producing
$2 \times 16 \times 256 \times 2 = 16384$~bytes/token/layer
(16.00\,KB), while the 10~full-attention layers use 4~KV heads
with head dimension~512, producing
$2 \times 4 \times 512 \times 2 = 8192$~bytes/token/layer
(8.00\,KB).
The weighted average across 60~layers is
$(50 \times 16.00 + 10 \times 8.00)/60 = 14.67$\,KB/token/layer,
an exact match with the measured density.

\begin{table}[t]
  \centering
  \small
  \caption{Analytical validation: theoretical vs.\ measured
    KV-cache density.  Theory uses
    $\rho = 2 H_\mathrm{kv} D E / 1024$~KB/tok/layer.}
  \label{tab:analytical}
  \begin{tabular}{@{}lrrrrrr@{}}
    \toprule
    \textbf{Model} & $H_\mathrm{kv}$ & $D$ & \textbf{dtype} & \textbf{Theory} & \textbf{Measured} & \textbf{Error} \\
    \midrule
    Pythia-1.4B  & 16 & 128 & bf16 &  8.00 &  8.00 &  0.0\% \\
    Gemma~4      & \multicolumn{2}{c}{\textit{weighted avg}} & bf16 & 14.67 & 14.67 &  0.0\% \\
    Qwen~3.6     &  2 & 256 & bf16 &  2.00 &  2.00 &  0.0\% \\
    gpt-oss      &  8 &  64 & bf16 &  2.00 &  1.12 & $-$44\% \\
    \bottomrule
  \end{tabular}
  \vspace{2pt}
  {\footnotesize Gemma~4 theoretical: $(50\times16.00 + 10\times8.00)/60$.
  gpt-oss gap: sliding-window layers store $\ll S$ tokens.}
\end{table}

\Cref{tab:analytical} shows three of four models match theory
exactly; the gpt-oss gap is structural.
Its sliding-window layers cap KV at a fixed window ($\approx$128
tokens) rather than the full $S=2048$ prefix, reducing their
effective density from 2.00 to $\approx$0.12\,KB/tok/layer.
The median across 36~alternating layers (18~full at 2.00 +
18~sliding at 0.12) yields $\approx$1.06, consistent with the
measured 1.12.

\subsection{Layer heterogeneity}
\label{sec:results:cv}

The layer coefficient of variation (CV) separates uniform from
hybrid architectures.
\Cref{tab:cv} presents the CV values.

\begin{table}[t]
  \centering
  \small
  \caption{Per-layer KV coefficient of variation.}
  \label{tab:cv}
  \begin{tabular}{@{}lrl@{}}
    \toprule
    \textbf{Model} & \textbf{CV} & \textbf{Pattern} \\
    \midrule
    Pythia-1.4B  & 0.000 & Uniform (all layers identical) \\
    Qwen~3.6     & 0.000 & Uniform (10 active layers, identical) \\
    Gemma~4      & 0.203 & Bimodal (local/global split) \\
    gpt-oss      & 0.902 & Extreme (alternating sliding/full) \\
    \bottomrule
  \end{tabular}
\end{table}

gpt-oss exhibits the highest layer CV in the KVScope dataset
at 0.902.
Its 36~profiled layers alternate between small sliding-window
layers (0.248\,MB each) and large full-attention layers
(4.211\,MB each), a 17$\times$ ratio.
This extreme heterogeneity has practical implications for KV
cache management: a paged allocator that assumes uniform per-layer
slot sizes will either waste memory on small layers or fail to
allocate contiguous blocks for large layers.

Gemma~4 shows moderate heterogeneity (CV\,=\,0.203) from its
50~sliding/10~full-attention split.
The sliding layers use 16~KV heads with head dimension~256
(16.00~KB/tok/layer), while the 10~full-attention layers use
4~KV heads with head dimension~512 (8.00~KB/tok/layer).
The weighted average $(50 \times 16.00 + 10 \times 8.00)/60 = 14.67$
exactly matches the measured density.

Pythia and Qwen~3.6 show CV near zero, confirming that all active
attention layers in these models have identical KV geometry.
Qwen~3.6's 10 active layers are uniform despite the model having
40~total layers; the 30~DeltaNet layers are invisible to the KV
profiler.

\subsection{KV footprint at 2048 tokens}
\label{sec:results:footprint}

\Cref{tab:footprint} presents absolute KV footprints at the
2048-token budget.

\begin{table}[t]
  \centering
  \small
  \caption{Median KV footprint at 2048 tokens (1024 for Pythia).}
  \label{tab:footprint}
  \begin{tabular}{@{}lrr@{}}
    \toprule
    \textbf{Model} & \textbf{KV MB} & \textbf{\% of 140\,GB VRAM} \\
    \midrule
    Qwen~3.6     & 42.4   & 0.03\% \\
    gpt-oss      & 79.9   & 0.06\% \\
    Pythia-1.4B  & 210.2  & 0.15\rlap{$^\dagger$} \\
    Gemma~4      & 1892.3 & 1.32\% \\
    \bottomrule
  \end{tabular}
  \vspace{2pt}
  {\footnotesize $^\dagger$At 1024 tokens; extrapolated to
  2048: $\approx$420\,MB (0.30\%).}
\end{table}

Gemma~4 dominates the KV footprint at 1.89\,GB for 2048~tokens,
consuming 1.32\% of the H200's 140\,GB VRAM from KV cache alone.
At longer sequences this grows linearly: a 32\,768-token context
would require approximately 30\,GB of KV cache for Gemma~4 alone,
leaving insufficient room for model weights on a single H200.

Hybrid SSM models are dramatically more efficient:
Qwen~3.6 at 42.4\,MB and gpt-oss at 79.9\,MB for 2048~tokens
represent 22--45$\times$ less KV memory than Gemma~4.
This efficiency comes from two sources:
(a)~fewer active KV layers (10 for Qwen~3.6 vs.\ 60 for Gemma~4)
and (b)~smaller KV head geometry (head dimension~64 for gpt-oss
vs.\ 256 for Gemma~4).

\subsection{Speculative decoding KV overhead}
\label{sec:results:spec}

The speculative decoding experiment pairs a 4-bit NF4-quantised
Llama-3.1-70B verifier with a bf16 Llama-3.2-1B draft model.
Of 20~prompts, 8~generated more than 10~tokens; the remaining
12~hit EOS immediately, a known issue with base (non-instruct)
models on short or adversarial prompts.

For the 8~valid prompts, median NVML overhead is 1804.6\,MB at
2048~tokens.
This includes KV caches for both models plus the 4-bit verifier
weights ($\approx$35\,GB) and draft model weights
($\approx$2\,GB).
The dual-model NVML overhead is approximately 1.8$\times$ that
of a single-model baseline at the same sequence length.
Since the NVML measurement includes model weights for both
models, the isolated KV-cache overhead ratio requires
per-layer hook data not available in the current tracer;
this is noted as a limitation.

Median throughput for valid prompts is 55.5\,tok/s, compared to
14.8\,tok/s for Gemma~4 and 19.1\,tok/s for Qwen~3.6.
The higher throughput reflects the 4-bit verifier's reduced
compute and the draft model's speculative token proposals.

\subsection{Post-EOS retention}
\label{sec:results:posteos}

The post-EOS score (\cref{eq:posteos}) measures what fraction of
generation-phase memory growth persists after the EOS token.
Pythia shows the highest retention: all 20~prompts have scores
between 0.067 and 0.888, with a median of~0.556.
This is consistent with PyTorch's CUDA allocator retaining freed
blocks in its caching pool rather than returning them to the
driver.
Gemma~4 shows moderate retention (median 0.330), while gpt-oss
and Nemotron-H show near-zero retention, suggesting more complete
cache release in hybrid architectures.
The speculative decoding runs show zero post-EOS retention,
likely because the 4-bit quantisation path uses a different
memory allocation strategy.

\subsection{Long-context KV growth validation}
\label{sec:results:longctx}

To validate that KV cache grows linearly with sequence length---as
predicted by \cref{eq:kvcache}---we sweep Pythia-1.4B from 128
to 8192~tokens on an NVIDIA RTX~4060~Laptop GPU (8\,GB VRAM).
\Cref{tab:longctx} reports NVML overhead at each generation length.

\begin{table}[t]
  \centering
  \small
  \caption{Long-context KV sweep for Pythia-1.4B (RTX~4060).
    Linear fit: $\mathrm{KV} = 0.1908\,S + 10.7$~MB,
    $R^2 = 0.99998$.}
  \label{tab:longctx}
  \begin{tabular}{@{}rrrr@{}}
    \toprule
    \textbf{Tokens} & \textbf{KV MB} & \textbf{KB/tok} & \textbf{KB/tok/lay} \\
    \midrule
    128   &   39.9 & 312 & 13.0 \\
    256   &   58.6 & 229 & 9.5 \\
    512   &  109.7 & 214 & 8.9 \\
    1024  &  202.9 & 198 & 8.3 \\
    2048  &  399.0 & 195 & 8.1 \\
    4096  &  791.0 & 193 & 8.1 \\
    8192  & 1575.1 & 192 & 8.0 \\
    \bottomrule
  \end{tabular}
\end{table}

The linear fit achieves $R^2 = 0.99998$, confirming O($S$) growth.
The measured slope of $0.1908$~MB/token (195.4~KB/token) is within
1.8\% of the theoretical $0.1875$~MB/token
($24 \times 8.00\,\text{KB} = 192\,\text{KB}$).
The 1.8\% excess is attributable to NVML allocator overhead beyond
the live KV tensors.

As sequence length increases, the per-token density converges to
the theoretical 8.00\,KB/tok/layer: at 128~tokens the density is
inflated by fixed allocator overhead, but by 8192~tokens it is
8.0\,KB/tok/layer.
This confirms that the density metric is asymptotically exact.

\paragraph{Cross-GPU consistency.}
The H200 measured KV tensor bytes for Pythia at 1024~tokens via
forward hooks: 210.2\,MB.
The RTX~4060 measured NVML overhead at 1024~tokens: 202.9\,MB.
The ratio is 1.04$\times$, confirming that KV-cache density is a
hardware-independent, architecture-determined constant.
The small difference arises because NVML on the H200 includes
PyTorch caching-allocator overhead while hook-based tensor byte
counts are exact.

\subsection{Failed and partial model attempts}
\label{sec:results:failures}

Several models could not be profiled successfully:
\begin{itemize}\setlength\itemsep{1pt}
  \item \textbf{LFM2.5-350M}: The \code{causal-conv1d} CUDA
    kernel had a symbol mismatch with PyTorch~2.12
    (\code{undefined symbol:
    \_ZN3c104cuda29c10\_cuda\_check\_implementationEiPKcS2\_jb}),
    preventing model loading.
    The V1 artifact had 19/20 prompts terminating at 0--9~steps
    due to missing chat template.
    LFM2.5 is excluded from V2 results.
  \item \textbf{Llama-3.1-70B standalone}: The bf16 model
    (140\,GB) exceeds the H200's usable VRAM, causing
    generation to stall indefinitely.
    4-bit NF4 quantisation resolves the VRAM issue but was only
    applied in the speculative decoding experiment.
    A standalone 4-bit Llama profile is not available.
  \item \textbf{GLM-4.7-Flash}: Failed during float8 weight
    conversion with
    \code{cat\_cuda not implemented for Float8\_e4m3fn}.
  \item \textbf{gpt-oss 4096-token run}: Crashed in the
    Triton/MXFP4 expert-parallel path due to a missing
    \code{device\_props.shared\_memory\_per\_block\_optin}
    attribute.
\end{itemize}

% =============================================================
\section{Data Quality Assessment}
\label{sec:quality}

\begin{table*}[t]
  \centering
  \small
  \caption{Artifact provenance and data-quality assessment.
    All V2 artifacts use a 2048-token budget (1024 for Pythia)
    and 20~prompts from the same prompt set.}
  \label{tab:data-quality}
  \resizebox{\textwidth}{!}{%
  \begin{tabular}{@{}p{0.14\linewidth}p{0.28\linewidth}p{0.42\linewidth}c@{}}
    \toprule
    \textbf{Model} & \textbf{Provenance} & \textbf{Notes} & \textbf{Grade} \\
    \midrule
    Pythia-1.4B & V2 run, 1024-token cap, NVML + hook data. & Clean KV traces; 24/24 layers uniform; post-EOS retention is allocator artefact. & A \\
    Gemma~4 & V2 rerun with chat template fix, 2048 tokens. & Strong KV traces; 60/60 layers profiled; 50 sliding + 10 full bimodal pattern confirmed. & A \\
    Qwen~3.6 & V2 run, 2048 tokens, 10/40 active KV layers. & Valid KV traces; DeltaNet layers invisible to profiler (expected). & A \\
    gpt-oss & V2 run, 2048 tokens, 36 layers. & Valid KV traces; extreme CV=0.902 from alternating layer sizes. & A \\
    Nemotron-H & V2 run, \code{use\_cache=False}, 2048 tokens. & No KV snapshots (SSM control); NVML overhead reflects model weights only. & D \\
    Spec Decode & V2 run, 4-bit verifier, 2048 tokens. & 8/20 valid prompts; no per-layer KV data in tracer; NVML overhead only. & B \\
    \bottomrule
  \end{tabular}}
\end{table*}

We assess each artifact on four criteria:
(1)~\emph{trace integrity}---all 20~prompts have valid KV
snapshots;
(2)~\emph{hardware provenance}---GPU identity is confirmed by
captured environment logs;
(3)~\emph{protocol uniformity}---token budget is 2048 across
all prompts within the artifact;
(4)~\emph{detector coverage}---per-layer KV byte counts are
available for density computation.
Overall grades: A~(all four criteria met), B~(partial coverage),
D~(control run, no KV traces).

The V2 dataset represents a significant improvement over V1:
all five transformer models use a 2048-token budget (1024 for
Pythia), environment metadata is fully captured, and the chat
template fix for Gemma~4 ensures full-length generation on all
20~prompts.

% =============================================================
\section{Discussion}
\label{sec:discussion}

\paragraph{What replicates.}
The cache-density hierarchy---local/global hybrid (14.67) $>$
pure MHA (8.00) $>$ dense GQA (4.00) $>$ hybrid SSM+attention
(1.12--2.00), all in KB/token/layer---is
consistent with theoretical predictions from \cref{eq:kvcache}
and with the H100 study~\cite{kvscope2026h100}.
The layer-CV signature of hybrid architectures replicates across
GPU generations: gpt-oss shows CV~0.902 (alternating layers),
Gemma~4 shows CV~0.203 (local/global split), and all
pure-attention or dense-GQA models show CV near zero.
These results confirm that KV-cache density is a stable,
architecture-determined fingerprint.

\paragraph{What is new in V2.}
Three findings are new:
(1)~gpt-oss's extreme layer heterogeneity (CV\,=\,0.902) was
masked in V1's 300-token runs; the V2 2048-token profile reveals
a 17$\times$ ratio between sliding-window and full-attention
layers.
(2)~Qwen~3.6's DeltaNet architecture leaves only 10/40 layers
with active KV tensors, producing a density of 2.00\,KB/tok/layer
despite the model's 35B parameter count.
(3)~Speculative decoding with a 70B verifier adds approximately
1.8$\times$ NVML overhead compared to single-model serving.

\paragraph{Practical implications.}
The 13$\times$ density range (1.12--14.67\,KB/tok/layer) means
that a serving system configured for one architecture will
dramatically misallocate for another.
A system tuned for GQA at 4\,KB/tok/layer would under-provision
by 3.7$\times$ for Gemma~4 and over-provision by 3.6$\times$
for gpt-oss.
The layer heterogeneity in gpt-oss (CV\,=\,0.902) further
complicates paged allocation: uniform slot sizes waste 47\% of
allocated memory on small layers.
Concretely, a uniform allocator must size every slot to the
largest layer (4.211\,MB for gpt-oss full-attention layers).
The 18~sliding-window layers each use only 0.248\,MB,
so total allocation is $36 \times 4.211 = 151.6$\,MB
versus actual usage of
$18 \times 4.211 + 18 \times 0.248 = 80.3$\,MB,
yielding $1 - 80.3/151.6 = 47\%$ waste.

\paragraph{The generation-length confound is resolved.}
V1 suffered from token budgets ranging 300--4096 across
artifacts, making direct footprint comparisons invalid.
V2's 2048-token budget (1024 for Pythia) eliminates this confound.
The density metric (KB/token/layer) is invariant to generation
length and remains the recommended comparison unit, but absolute
footprint comparisons are now also valid within V2.

% =============================================================
\section{Limitations}
\label{sec:limitations}

\begin{enumerate}\setlength\itemsep{1pt}
  \item Pythia-1.4B uses a 1024-token cap (not 2048) due to
    its shorter generation horizon; absolute footprint
    comparisons with Pythia require extrapolation.
  \item The speculative decoding experiment has only 8/20 valid
    prompts (base model EOS sensitivity) and lacks per-layer KV
    data in its tracer output.
  \item LFM2.5-350M is excluded due to a CUDA kernel
    incompatibility; coverage of pure linear-attention models
    remains a gap.
  \item 4-bit NF4 quantisation is used for the 70B verifier in
    speculative decoding; KV tensor byte counts may differ
    slightly from bf16 due to quantisation-aware allocation.
  \item Nemotron-H's NVML overhead (132\,GB) reflects model
    weights, not KV cache; this artifact serves only as a
    zero-KV control.
  \item Model loading failures for GLM-4.7-Flash and the
    gpt-oss long-run limit coverage of very large or
    non-standard models.
\end{enumerate}

% =============================================================
\section{Conclusion}
\label{sec:conclusion}

KVScope~V2 on H200 confirms that architecture-level KV-cache
density is a stable, reproducible fingerprint across GPU
generations, now with a standardised 2048-token protocol
(1024 for Pythia) that enables direct cross-model comparisons.
The density hierarchy spans 13$\times$:
gpt-oss hybrid SSM+attention at 1.12\,KB/token/layer,
Qwen~3.6 DeltaNet+GQA at 2.00,
dense GQA at 4.00,
pure MHA at 8.00,
and Gemma~4 local/global hybrid at 14.67.
Layer heterogeneity ranges from CV\,=\,0 (uniform architectures)
to CV\,=\,0.902 (gpt-oss alternating layers), with direct
implications for paged KV-cache allocation.
The SSM control (Nemotron-H) confirms O(1) memory growth,
and speculative decoding adds 1.8$\times$ NVML overhead.

Practical failures for LFM2.5, GLM-4.7-Flash, and the gpt-oss
long-run highlight that CUDA kernel compatibility and
expert-parallel Triton paths require explicit GPU-specific
support not yet present in general-purpose profiling frameworks.
The strongest recommendation for a V3 study is to extend coverage
to pure linear-attention models (once kernel compatibility is
resolved), add long-context sweep data to characterise
non-linear KV growth, and integrate H2O eviction profiling to
quantify the memory--quality tradeoff analytically.

% =============================================================
\section*{Artifact Availability}
\label{sec:artifacts}

All metrics in this paper are extracted from JSON files under
\texttt{results\_v2/} by the scripts in \texttt{scripts/}.
The artifact package includes: all profile JSONs, model config
files, \code{pip freeze} output, full \code{nvidia-smi -q}
output, CPU/OS details, and 26~run logs.
The V2 artifact is archived at Zenodo (DOI to be assigned upon
acceptance).
The H100 V1 artifact is archived at
\url{https://doi.org/10.5281/zenodo.19871039}~\cite{kvscope2026h100}.
Code and reproduction scripts are available at
\url{https://github.com/CosmicAlgo/kv-cache-dynamics}.

\printbibliography
\end{document}
